\documentclass[12pt,final,twoside]{report}
%%%%%%%%%%%%%%%%%%%%%%%%%%%%%%%%%%%%%%%%%%%%%%%%%%%%%%%%%%%%%
% Meta informations:
\newcommand{\trauthor}{Vincent Dahmen}
\newcommand{\trmatrikelnummer}{6689845}
\newcommand{\trtype}{Bachelorthesis} %{Diplomarbeit} %{Dissertation}
\newcommand{\trtitle}{Evaluation of prefix constraints and their impact on related and unrelated language pairs}
\newcommand{\tremail}{4dahmen@informatik.uni-hamburg.de}
\newcommand{\trinstitute}{Dept. Informatik -- Knowledge Technology, WTM}
\newcommand{\trwebsiteordate}{{https://www.informatik.uni-hamburg.de/WTM/}}
\newcommand{\trdate}{\today}

\newcommand{\trcourseofstudies}{Informatik} %{Bioinformatik} %{Intelligent Adaptive Systems} 
\newcommand{\trauthordegree}{} %{\ B.Sc.}
\newcommand{\trgutachterA}{\href{mailto:alpay@informatik.uni-hamburg.de}{M. Sc. Tayfun Alpay}}
\newcommand{\trgutachterB}{\href{mailto:heinrich@informatik.uni-hamburg.de}{Dr. Stefan Heinrich}}
\newcommand{\trfach}{Knowledge Technology, WTM}
\newcommand{\trkeywords}{Learning, Language, Demonstration}
%%%%%%%%%%%%%%%%%%%%%%%%%%%%%%%%%%%%%%%%%%%%%%%%%%%%%%%%%%%%%
% Languages:

% Falls die Ausarbeitung in Deutsch erfolgt:
% \usepackage[german]{babel}
% \usepackage[T1]{fontenc}
% \usepackage[latin1]{inputenc}
% \usepackage[latin9]{inputenc}                     
% \selectlanguage{german}

% If the thesis is written in English:
\usepackage[english]{babel}                         
\selectlanguage{english}

%%%%%%%%%%%%%%%%%%%%%%%%%%%%%%%%%%%%%%%%%%%%%%%%%%%%%%%%%%%%%
% Bind packages:
\usepackage{acronym}                    % Acronyms
\usepackage{algorithmic}                % Algorithms and Pseudocode
\usepackage{algorithm}                  % Algorithms and Pseudocode
\usepackage{amsfonts}                   % AMS Math Packet (Fonts)
\usepackage{amsmath}                    % AMS Math Packet
\usepackage{amssymb}                    % Additional mathematical symbols
\usepackage{amsthm}
\usepackage{booktabs}                   % Nicer tables
%\usepackage[font=small,labelfont=bf]{caption} % Numbered captions for figures
\usepackage{color}                      % Enables defining of colours via \definecolor
\definecolor{uhhRed}{RGB}{226,0,26}     % Official Uni Hamburg Red
\definecolor{uhhGrey}{RGB}{136,136,136} % Official Uni Hamburg Grey
\definecolor{uhhLightGrey}{RGB}{220, 220, 220}
\usepackage{fancybox}                   % Put equations in a frame
\usepackage{fancyhdr}                   % Packet for nicer headers
%\usepackage{fancyheadings}             % Nicer numbering of headlines
\usepackage[body={5.8in,9in}]{geometry} % Type area (size, margins...)

%\geometry{a4paper,outer=3.35cm}        % !!!Release version (Normal margins)
%\geometry{a4paper,outer=2.5cm}         % !!!Print version (Additional margin on the left for the binding)
%\geometry{a4paper}                     % !!!Proofread version (Additional margin on the right for corrections)
\geometry{a4paper,outer=3.134cm}        % !!!Draft version (Same margins on left and right)
%\geometry{paperheight=10.0in,paperwidth=6.4in,top=0.51in,left=0.3in}  % !!!Developer version (Minimal margins)

\usepackage{graphicx}                   % Inclusion of graphics
%\usepackage{latexsym}                  % Special symbols
\usepackage{longtable}                  % Allow tables over several parges
\usepackage{listings}                   % Nicer source code listings
\usepackage{multicol}                   % Content of a table over several columns
\usepackage{multirow}                   % Content of a table over several rows
\usepackage{rotating}                   % Alows to rotate text and objects
%\usepackage{siunitx}					% Nicer display of units and number according to SI standards
%\sisetup{tight-spacing=true}  
\usepackage[hang]{subfigure}            % Allows to use multiple (partial) figures in a fig
%HINT: subfigure may be depricated already - maybe try subfig instead 
%\usepackage[font=footnotesize,labelfont=rm]{subfig}    % Pictures in a floating environment
\usepackage{tabularx}                   % Tables with fixed width but variable rows
\usepackage{url,xspace,boxedminipage}   % Accurate display of URLs

%%%%%%%%%%%%%%%%%%%%%%%%%%%%%%%%%%%%%%%%%%%%%%%%%%%%%%%%%%%%%
% PDF Information und Definitions:
\author{\trauthor}

\ifx\pdftexversion\undefined
\usepackage{hyperref}
\else
\usepackage[colorlinks=false,           % link is colores (true) or has colored frame (false)
            linkcolor=uhhRed,           % case colorlinks=true: define color.
            urlcolor=uhhRed,
            citecolor=uhhRed,
            bookmarks,                  % Place bookmarks erstellen
            bookmarksopen=true,         % Bookmarks will be shown at start (true/false)
            pdfpagemode=UseOutlines,    
            bookmarksopenlevel=1,       % Define the depth of shown links
            bookmarksnumbered,          % Numbers of chapers in Bookmarks
            pdftitle={\trtitle},
            pdfsubject={\trtype},
            pdfkeywords={\trkeywords},
            pdfauthor={\trauthor},
            plainpages=false
            ]{hyperref}
\fi

\usepackage{apacite}			% Extended cite options/style; Must be loaded after hyperref

\ifx\pdftexversion\undefined
\else
\pdfoutput=1                            % Disable PDF-Output
\pdfimageresolution=1200
\pdfcompresslevel=2                     % 0 = no compression, 9 = strongest compression
\fi

%%%%%%%%%%%%%%%%%%%%%%%%%%%%%%%%%%%%%%%%%%%%%%%%%%%%%%%%%%%%%
% Configurationen:

\hyphenation{whe-ther}                  % Manually use: "\-" in a word: Staats\-ver\-trag

%\lstloadlanguages{C}                   % Set the default language for listings
\DeclareGraphicsExtensions{.pdf,.svg,.jpg,.png,.eps} % first try pdf, then eps, png and jpg
\graphicspath{{./img/}}                 % Path to a folder where all pictures are located
\pagestyle{fancy}                       % Use nicer header and footer

% Redefine the environments for floating objects:
\setcounter{topnumber}{3}
\setcounter{bottomnumber}{2}
\setcounter{totalnumber}{4}
\renewcommand{\topfraction}{0.9}        %Standard: 0.7
\renewcommand{\bottomfraction}{0.5}     %Standard: 0.3
\renewcommand{\textfraction}{0.1}       %Standard: 0.2
\renewcommand{\floatpagefraction}{0.8}  %Standard: 0.5

% Tables with a nicer padding:
\renewcommand{\arraystretch}{1.2}

% Chapter and Sections will not be written in capitals
\renewcommand{\chaptermark}[1]{\markboth{\chaptername.\ \thechapter.\ #1}{}}
\renewcommand{\sectionmark}[1]{\markright{\thesection.\ #1}}

% Avoid french spacing (double spacing after a full stop)
\frenchspacing

%%%%%%%%%%%%%%%%%%%%%%%%%%%%
% Additional 'theorem' and 'definition' blocks:
\theoremstyle{plain}
\newtheorem{theorem}{Theorem}[chapter]
%\newtheorem{theorem}{Satz}[chapter]    % Wenn in Deutsch geschrieben wird.
\newtheorem{axiom}{Axiom}[chapter]     
%\newtheorem{axiom}{Fakt}[chapter]      % Wenn in Deutsch geschrieben wird.
%Usage:%\begin{axiom}[optional description]%Main part%\end{fakt}

\theoremstyle{definition}
\newtheorem{definition}{Definition}[chapter]

%Additional types of axioms:
\newtheorem{lemma}[axiom]{Lemma}
\newtheorem{observation}[axiom]{Observation}

%Additional types of definitions:
\theoremstyle{remark}
%\newtheorem{remark}[definition]{Bemerkung} % Wenn in Deutsch geschrieben wird.
\newtheorem{remark}[definition]{Remark} 

%%%%%%%%%%%%%%%%%%%%%%%%%%%%
% Provides TODOs within the margin:
\newcommand{\TODO}[1]{\marginpar{\emph{\small{{\bf TODO: } #1}}}}

%%%%%%%%%%%%%%%%%%%%%%%%%%%%
% Abbreviations and mathematical symbols
\newcommand{\modd}{\text{ mod }}
\newcommand{\RS}{\mathbb{R}}
\newcommand{\NS}{\mathbb{N}}
\newcommand{\ZS}{\mathbb{Z}}
\newcommand{\dnormal}{\mathit{N}}
\newcommand{\duniform}{\mathit{U}}

\newcommand{\erdos}{Erd\H{o}s}
\newcommand{\renyi}{-R\'{e}nyi}
% it is recommented to define complex terms as expression/newcommand and use the expression in the tex instead.

%%%%%%%%%%%%%%%%%%%%%%%%%%%%%%%%%%%%%%%%%%%%%%%%%%%%%%%%%%%%%
% Document:

\begin{document}

\pagenumbering{Roman}                   % Roman pagenumbering for lists and meta pages
\renewcommand{\headheight}{14.5pt}      % Size of headings

\thispagestyle{empty}
\fancyhead[LO,RE]{}                     % Define the header style for the meta pages

%%%%%%%%%%%%%%%%%%%%%%%%%%%%
% Cover sheet

\begin{titlepage}
%---Possibility 1:
    \begin{flushleft}
        \includegraphics[width=67mm]{uhhLogoL.pdf}\\
    \end{flushleft}
    \rule{\textwidth}{0.4pt}
        \newline
        \vspace{2.0cm}
        \begin{center}
          \LARGE \textbf{\trtitle}
        \end{center}
    \vspace{2.0cm}
    \begin{center}
      \textbf{\trtype}\\
      %im Arbeitsbereich \trfach\\
      at Research Group \trfach\\
      \trgutachterA\medskip\\
      %Fachbereich Informatik\\
      Department of Informatics\\
      %MIN-Fakult\"at\\
      MIN-Faculty\\
      Universit\"at Hamburg \\[1.0cm] %not allowed to translate Universit{\"a}t Hamburg
      %vorgelegt von \\
      submitted by \\
      \textbf{\href{mailto:\tremail}{\trauthor\trauthordegree}}\\
      %Studiengang:   \trcourseofstudies \\
      Course of study:   \trcourseofstudies \\
      Matrikelnr.:  \trmatrikelnummer \\
      %a\\
      on\\
      \trdate
    \end{center}
    \vspace{2cm}
    \begin{center}
    \begin{tabular}{ll}
    %Gutachter: & \trgutachterA \\
    Examiners: & \trgutachterA \\
                   & \trgutachterB \\
    %Betreuung: & \trbetreuung \\    	% Adviser are not allowed to demand getting mentioned here, but are happy getting credited by student's initiative
    \end{tabular}
    \end{center}
    \vfill
%    \begin{tabular}{l}
%    \trauthor \\
%    Matrikelnummer:  \trmatrikelnummer \\
%    \trstrasse \\
%    \trort
%    \end{tabular}
%    \newline
    \rule{\textwidth}{0.4pt}
    \newpage 
\end{titlepage}

    %backsite of cover sheet is empty!
\thispagestyle{empty}
\hspace{1cm}
\newpage

%%%%%%%%%%%%%%%%%%%%%%%%%%%%
% Abstract:
\section*{Abstract}\label{sec:abstract}
%\fancyhead[LE,RO]{\it Abstract}
%\addcontentsline{toc}{chapter}{\numberline{}Abstract}
\input{"abstract-en.tex"}

% If the abstract ist not longer than half a page, then the German Zusammenfassung can be places on the same page
%\cleardoublepage
\vfill
\section*{Zusammenfassung}\label{sec:zusammenfassung}
%\addcontentsline{toc}{chapter}{\numberline{}Zusammenfassung} %Add the Zusammenfassung to the TOC
\input{"abstract-de.tex"}
\fancyhead[LE,RO]{\it Abstract}

\cleardoublepage

%%%%%%%%%%%%%%%%%%%%%%%%%%%%
% Lists:
%\setcounter{tocdepth}{1}               % depth of the table of contents (for BSc and MSc Thesis 1 is recommented)
\fancyhead[LE,RO]{\it Contents}
\tableofcontents
\cleardoublepage
% List of Figures and List of tables are optional. -> Not needed in most theses.
%\fancyhead[LE,RO]{\it List of Figures}
%\listoffigures
%\cleardoublepage
%\fancyhead[LE,RO]{\it List of Tables}
%\listoftables
%\cleardoublepage
%\lstlistoflistings
%\cleardoublepage

\fancyhead[LE]{\it \leftmark}           % Define the header style for the text pages
\fancyhead[RO]{\it \rightmark}          % Define the header style for the text pages
\fancyhead[LO,RE]{}                     % Define the header style for the text pages

%%%%%%%%%%%%%%%%%%%%%%%%%%%%
% The content will be included here:
\pagenumbering{arabic}

\input{"chapters.tex"}


%\fancyhead[LE]{\it \leftmark}
%\chapter{}
\fancyhead[LE,RO]{\it Bibliography}       % A bibliography never have a letter or numbering!
    \bibliographystyle{apacite}             % Style for presenting the literature
	\phantomsection\addcontentsline{toc}{chapter}{Bibliography}% Add to the TOC
    \bibliography{bib}
\cleardoublepage

%%%%%%%%%%%%%%%%%%%%%%%%%%%%
% Formal page 1
\vspace{2cm}
\chapter*{Erkl\"arung der Urheberschaft}
\label{sec:urheber}
\fancyhead[LE]{\it Erkl\"arung der Urheberschaft}
	Hiermit versichere ich an Eides statt, dass ich die vorliegende
	\trtype{} im Studiengang \trcourseofstudies{}
	selbstst{\"a}ndig verfasst und keine anderen als die angegebenen
	Hilfsmittel - insbesondere keine im Quellenverzeichnis nicht
	benannten Internet-Quellen – benutzt habe. Alle Stellen, die
	w{\"o}rtlich oder sinngem{\"a}{\ss} aus Ver{\"o}ffentlichungen entnommen wurden,
	sind als solche kenntlich gemacht. Ich versichere weiterhin, dass
	ich die Arbeit vorher nicht in einem anderen Pr{\"u}fungsverfahren
	eingereicht habe und die eingereichte schriftliche Fassung der
	auf dem elektronischen Speichermedium entspricht.
	
\vspace{4cm}
\noindent Ort, Datum \hfill Unterschrift

%The backcover is always empty
\newpage
\thispagestyle{empty}
\hspace{1cm}
\newpage

%%%%%%%%%%%%%%%%%%%%%%%%%%%%
% Formal page 2
\vspace{2cm}
\chapter*{Erkl\"arung zur Ver\"offentlichung}
\label{sec:urheber}
\fancyhead[LE]{\it Erkl\"arung zur Ver\"offentlichung}
Ich stimme der Einstellung der \trtype{} in die Bibliothek des Fachbereichs Informatik zu.

\vspace{4cm}
\noindent Ort, Datum \hfill Unterschrift

%The backcover is always empty
\newpage
\thispagestyle{empty}
\hspace{1cm}
\newpage

\end{document}
%EOF

